\section{Introduction}
\label{sec:introduction}

Schizophrenia is a severe neuropsychiatric disorder affecting approximately 1\% of the global population, characterized by deficits in cognitive control, social cognition, and visual attention \cite{holzman1973}. Clinical diagnosis remains largely qualitative and subject to inter-rater variability, motivating the search for objective, quantifiable biomarkers. Atypical eye movement patterns during free-viewing paradigms have emerged as one of the most replicated neurobiological markers of schizophrenia \cite{benson2012}: patients exhibit restricted visual exploration, social gaze avoidance, and abnormal saccadic dynamics compared to healthy controls \cite{morita2020}.

Existing approaches to eye movement-based schizophrenia classification fall into two broad categories. Hand-crafted feature methods \cite{shpigelman2008, huang2020, tseng2013} extract static summary statistics over fixation sequences and classify them with shallow models, discarding the spatiotemporal structure of scanpaths. Recent deep learning methods such as MSNet \cite{song2025} achieve state-of-the-art performance by aggregating fixation density maps into multi-scale CNN representations; however, this collapses the temporal saccade trajectory and yields no fixation-level clinical interpretability.

To address these limitations, we propose a \textbf{Hierarchical 5-Tier Framework} that systematically integrates clinical domain knowledge with deep representation learning. Our contributions are threefold: 
(1) we introduce Tier 2 contextual delta features (60 cross-category contrast features) that encode stimulus-specific visual response differences clinically linked to social cognition deficits; 
(2) we develop GNN-CEFAM, a spatiotemporal Graph Attention Network featuring a Cross-Attention Enhanced Fusion Module (CEFAM) that computes genuine, non-degenerate cross-attention over pre-pooling fixation nodes; and 
(3) we propose BiCA-HS, a Transformer-based hybrid stream model that fuses raw fixation trajectories with a 135-dimensional clinical expert vector in a unified attention space. Evaluated across three random seeds, our Tier 5 logistic regression meta-learner achieves a state-of-the-art out-of-fold AUC of $0.9590 \pm 0.0009$, achieving higher diagnostic performance than the published MSNet test benchmark (reported AUC\,=\,0.8854).
